\thispagestyle{empty}

\begin{center}
  {\large\bfseries \ProjectTitleEng}\\
\end{center}

\begin{center}
  \AuthorOne\\
  \AuthorTwo\\
  \vspace{0.7cm}
  \noindent{\textbf{Keywords}:}\\
  Implied volatility, IBEX options, Black-76, machine learning,
  explainable AI, SHAP, dashboard, API, cloud\\
  \vspace{0.7cm}
  \noindent{\textbf{Abstract}}\\
\end{center}

The main objective of this Master's Thesis is to build implied-volatility models for IBEX options that are as reliable and powerful as modern artificial intelligence models, while remaining as explainable as traditional mathematical models used in quantitative finance. The work addresses a central tension: AI models can capture complex nonlinear relationships across the volatility surface, whereas classical models provide structure, interpretability and auditability. The proposed solution aims to combine both properties in a practical modelling framework.

Starting from real trades, contracts, underlying futures and EONIA/\texteuro{}STR rate series, the project builds an implied-volatility database through financial validations and numerical inversion of the Black-76 pricing model. Several regression families are then trained, including linear regression, Random Forest, XGBoost, neural networks and a quantum-inspired architecture, based on quantum computing concepts as its name suggests. The evaluation is not restricted to predictive error; it also considers temporal stability, overfitting control, financial coherence of the surface, and regional behaviour across moneyness and maturity.

The main contribution is to turn predictive models into explainable, auditable and queryable artefacts through a dashboard and a FastAPI service. Explainability is addressed through global and local SHAP, volatility surfaces, ICE and ALE curves, surrogate trees, symbolic regression, historical neighbours and residual diagnostics. This combination makes it possible to justify predictions in a language close to traditional financial models, without giving up the approximation power of AI. The resulting architecture is reproducible and oriented towards model validation, internal audit and operational use in financial environments.

Additionally, an API has been enabled so that anyone can run and use these models.

\newpage
